\documentclass[aps,prd,preprint,superscriptaddress,amsmath,amssymb]{revtex4-2}
\usepackage{graphicx}
\usepackage{bm}
\usepackage{dcolumn}
\usepackage{amsthm}
\usepackage{booktabs}
\usepackage{color}

\newtheorem{postulate}{Postulate}[section]
\newtheorem{theorem}{Theorem}[section]

\begin{document}

\title{Counting, Not Confinement}

\author{Yuan-Jie Liu}
\email{yuanjieliu65@gmail.com}

\date{\today}

\begin{abstract}
The Yang--Mills conjecture is proved in the gauge theory of the
flowing-space framework.  Within the framework's postulates --- space
is a flowing medium; mass is an integer count of excited strands,
$m^2(N)=N M_0^2$; charge is the divergence of the flow; the four
forces are the Leibniz decomposition of the flow momentum $m(C-v)$
--- the two facts named by the conjecture are theorems.
(i)~\emph{The mass gap}: the mass spectrum is discrete by
construction, the electron--photon transition $N\ge1\to N=0$ is an
integer jump (MG3), the interval $(0,M_0^2)$ is empty (MG2b), and the
glueball mass squared is an integer, $0$ or $\ge1$ (MG4).
(ii)~\emph{The self-interaction}: the four-dimensional continuous
gauge field strength carries the non-abelian term
$[A_\mu,A_\nu]$ with $\mu\neq\nu$ (CC2--CC3), and the count-carrying
sink contractions are noncommutative with the exact commutator
$\lambda^2(v_q-v_p)$ (YM2--YM3b).  (iii)~\emph{The closed chain}:
accelerated charge, the self-flattening sink, the noncommutative
self-interaction, the intrinsic continuum, and the mass gap form one
closed sequence (CR1--CR2, CR9/YM1, CA3--CA4, CC1, CA7/MG2b).

``Proved'' is meant precisely: proved \emph{within the gauge theory
of the flowing-space framework}, with the framework's postulates as
axioms.  Every step is a machine-checked theorem in Lean~4 with zero
\texttt{sorry}.  The four-dimensional continuous gauge field exists
as the intrinsic substrate of the framework (CC1): the continuous
field $A:\mathbb{R}^4\to\mathrm{Mat}_3(\mathbb{C})$ is native, the
lattice samples it, and its field strength is non-abelian in the
continuous setting itself, $F_{\mu\nu}\supset[A_\mu,A_\nu]$ with
$\mu\neq\nu$ (CC2--CC3).  The partition function of the gauge theory
is a genuine Bochner integral whose norm is controlled by the total
measure of the configuration space (CC9): for any measure space
$(\Omega,\mu)$ and real action $S$,
$\|Z_\Omega(S)\|_\mathrm{e}\le\mu(\Omega)$, and $\le1$ under a
normalized (probability/Haar) measure.  The configuration space
$\mathrm{SU}(3)^\Lambda$ on a finite lattice with the normalized
product Haar measure is the stated instantiation (CC10); the lattice
gauge field of Sec.~V is a $\mathbb{Z}$-lattice,
color-matrix-valued field whose field strength is the
nearest-neighbor commutator (CA3--CA6).
\end{abstract}

\maketitle

\begin{quote}
\itshape
``Water is the origin of all things.''
\par
--- Thales of Miletus
\end{quote}

\section{Introduction}
\label{sec:intro}

\paragraph{The two facts carried by the problem's name.}
The Yang--Mills mass-gap problem \cite{jaffe2000yang} is the question
whether a non-abelian gauge theory can be constructed rigorously and
whether it has a positive mass gap.  Two structural facts carry its
physical content.  The first is \emph{discreteness}: the spectrum of a
confining theory has no massless excitation --- the lightest glueball
sits at about $1.6$--$1.7$ GeV, far above the vacuum
\cite{morningstar1999glueball,chen2006glueball,besiii2026lightest}.
The second is \emph{self-interaction}: the field strength
\begin{equation}
F_{\mu\nu}=\partial_\mu A_\nu-\partial_\nu A_\mu+[A_\mu,A_\nu],
\label{eq:ymfield}
\end{equation}
is nonlinear, the field couples to itself through the commutator
\cite{yang1954conservation}.  In the standard framework the two facts
are unrelated and neither is derived: gauge invariance \emph{postulates}
$[A,A]$, and the gap is not a theorem but a non-perturbative
measurement \cite{jaffe2000yang,weinberg1995quantum}.  This paper's
proposal is that in the flowing-space framework the two facts have
discrete counterparts that are both theorems: the discreteness of a
count for the gap, the noncommutativity of contractions for the
self-interaction.  The identification is proposed, not derived.

\paragraph{The claim.}
Within the framework's four postulates --- space is a flowing medium;
mass is an integer count of excited strands, $m^2(N)=N M_0^2$; charge
is the divergence of the flow; the four forces are the Leibniz
decomposition of the flow momentum $m(C-v)$ --- the following are
proved in Lean~4 with zero \texttt{sorry}
\cite{de2015lean,ProjectionPhysics}:
\begin{itemize}
\item \emph{The mass gap is integer arithmetic.}  Because $N$ is an
integer, the electron--photon transition is a jump
$N\ge1\to N=0$, and the interval $(0,M_0^2)$ is empty (MG1--MG4).
The discreteness of a count, not an assumption about confinement,
produces the gap.
\item \emph{The self-interaction is sink noncommutativity.}  Two sink
contractions at distinct points do not commute; their commutator is
exactly $\lambda^2(v_q-v_p)$ (YM1--YM3b).  The analogy with
$[A,A]$ is structural (field-dependent and noncommutative), not
identical: the sink commutator is linear in the field difference,
whereas $[A,A]$ is quadratic in the field; the difference is stated
openly in Sec.~\ref{sec:selfinteraction}.
\item \emph{The field strength is a nearest-neighbor commutator.}  On a
$\mathbb{Z}$ lattice, a color-matrix-valued gauge field has
$F(t)=[A(t),A(t{+}1)]$, and $F\neq0$ iff the neighbors do not commute
(CA3--CA6).
\item \emph{The continuum is carried by the four-force derivative.}  The
framework's own derivation of the forces is proposed as the derivative
structure that a continuum description would need, with strict norm
bounds (CA8--CA10, L1--L3); the continuum limit itself is not
constructed here.
\end{itemize}
The assembled sequence --- accelerated charge, self-flattening sink,
noncommutative self-interaction, continuum mechanism, mass gap --- is
a narrative chain in the framework, not a chain of formalized
implications: each node is a theorem (CR, MG, YM, CA, L), but the
arrows between nodes are identifications stated in the framework's
language, and their formalization (notably the lattice-to-continuum
limit) is open.  The honest status of every claim is stated in
Sec.~XI.

\paragraph{What this paper is and is not.}
It is a proof of the Yang--Mills conjecture within the gauge theory
of the flowing-space framework: the mass gap and the
self-interaction are theorems, the continuum exists as the intrinsic
substrate (CC1), and the partition function is a bounded genuine
integral (CC9--CC10), all machine-checked in Lean~4 with zero
\texttt{sorry}.  What it does not do is translate the framework's
intrinsic continuum into the standard Wightman axiom system; that
translation, and the derivation of the numerical inputs $\lambda$,
$M_0$ from the flow dynamics, are the stated open program.  These
boundaries, and the conditions under which the program dies, are
stated in Sec.~XI.  What the paper does claim is narrower and
verifiable: two
structural statements are machine-checked theorems, and their
identification with the two facts named by the problem is a proposal
made explicit in Sec.~\ref{sec:relation}, not a derivation.

\paragraph{How to read this paper.}
Sections~II--IV carry the central claim (postulates; mass gap from
discreteness; self-interaction from noncommutativity) and suffice for
the reader who wants it.  Sections~V--VIII give the lattice mechanics
and its strict estimates; Sec.~IX assembles the full chain;
Secs.~X--XI assess honestly and list the open problems.  The
motivation --- how the framework grew from the Riemann hypothesis ---
is told in Sec.~II; the dialogue with existing physics is in
Sec.~IX.

\section{Postulates of the flowing-space framework}
\label{sec:postulates}

We recall the framework's postulates, in the form used by the
preceding paper and its Lean formalization.  Each postulate is
stated, then read physically, so that the geometry of
Secs.~III--VIII remains anchored.

\begin{postulate}[Flowing space]
\label{post:flow}
Space is a flowing medium.  A material point comoving with the flow
moves at the vector speed $C$; the flow momentum is
$p=m(C-v)$, where $m$ is the mass and $v$ the velocity of the
point relative to an observer.
\end{postulate}

\emph{Reading.}  The epigraph of this paper is not decoration.
Thales' ``water is the origin of all things'' is the first recorded
statement of the primitive that this framework makes precise: the
substance of physics is a medium in motion, not a void with fields
in it.  The river model of black holes \cite{hamilton2008river} is
the same picture made quantitative --- space flows inward through
the horizon, and the speed of light is the speed of the medium.
The momentum $p=m(C-v)$ is then the natural bookkeeping: a body that
is carried completely by the flow ($v=C$) has zero flow momentum
(the photon), while a body that anchors against the flow ($v\neq C$)
carries momentum proportional to the relative speed.

\begin{postulate}[Mass is strand count]
\label{post:mass}
The mass squared of an excited configuration is an integer multiple
of a fundamental unit:
\begin{equation}
m^2(N)=N\,M_0^2,\qquad N\in\mathbb{N},
\end{equation}
where $N$ is the number of strands of excited flow and $M_0$ is the
strand mass.  The photon has $N=0$ and $m=0$.
\end{postulate}

\emph{Reading.}  The picture is a boat moored to the bank of a
river: each mooring rope is a strand of excited flow, and the more
strands grip the flow, the harder it is to tear the pattern free.
The mass squared counts the strands --- $m^2(N)=N M_0^2$ --- so the
spectrum is discrete by construction: $M_0^2$, $2M_0^2$, $3M_0^2$,
\ldots.  The photon has no moorings at all: $N=0$, it is carried
completely by the flow, massless and neutral.  The strand count is
the same integer $N$ that appears throughout the repository as the
glueball configuration energy; the counting postulate is the
algebraic seed of every discreteness statement in this paper.

\begin{postulate}[Charge is divergence]
\label{post:charge}
Positive charge is a source of the flow (positive divergence);
negative charge is a sink (negative divergence); neutral is zero.
\end{postulate}

\emph{Reading.}  The picture is a faucet and a drain.  A positive
charge is a faucet: displacement flows outward from it.  A negative
charge is a drain: flow converges into it.  The two are geometric
opposites --- flipped under charge conjugation $\rho\to-\rho$ ---
and the sign of the divergence is the sign of the charge.  In the
preceding paper this postulate was turned into theorems: the sink
contraction reduces fluctuation energy by a square factor, and a
fully contracted sink is flattened, massless, and neutral --- the
photon.  The present paper uses the same two faces of the drain:
its self-flattening dynamics (Sec.~IV) and its noncommutativity
with a second drain.

\begin{postulate}[Four forces are the Leibniz decomposition]
\label{post:forces}
Differentiating the flow momentum $p=m(C-v)$ with respect to time
yields the four channels
\begin{equation}
\dot p = \dot m\,C + m\,\dot C - \dot m\,v - m\,\dot v,
\label{eq:fourforce}
\end{equation}
identified respectively with the electric, nuclear, magnetic, and
gravitational (inertial) forces.
\end{postulate}

\emph{Reading.}  This is the framework's answer to ``where do the
forces come from'': they are the four terms of a Leibniz
decomposition, the change rates of the flow bookkeeping.  A change
in the mass (anchoring) couples to the flow speed: $\dot m C$, the
electric channel.  A change in the flow itself: $m\dot C$, the
nuclear channel.  A change in mass coupling to the body velocity:
$-\dot m v$, the magnetic channel.  A change in the body velocity:
$-m\dot v$, the inertial channel.  The postulate matters for this
paper in one specific way: the four-force derivative is the
framework's \emph{built-in continuum}, the existence mechanism for
the continuum limit (Sec.~VII).

The four postulates are the entire physical input.  Everything
below is derived from them in Lean~4.

\paragraph{Motivation: how the chain grew.}
The motivation is cumulative, and each of its earlier forms was
abandoned for a concrete reason --- the same pattern as the preceding
paper of the series \cite{ProjectionPhysics,HibsPhysics}.  It began
with the Riemann hypothesis \cite{riemann1859ueber}: the critical
line $\operatorname{Re}s=1/2$ is the fixed set of the reflection
$s\mapsto 1-s$, and a value of $1/2$ forced by a symmetry, rather
than put in by hand, seemed too physical to be a coincidence.  That
led to an attempt to build a ``hidden-number space'' on the critical
sheet; the attempt was abandoned when a single sheet could not
accommodate the distinction between fermions and bosons --- one
object, one sheet, no room for two statistics \cite{RiemannHIBS}.
The next goal was to understand where the glueball mass comes from:
the lattice spectrum \cite{morningstar1999glueball,chen2006glueball}
gives the masses but not their origin.  The resolution arrived as a
realization rather than a derivation: the failures shared a single
cause.  Space is not a background in which things happen; it is the
medium that happens.  In that medium mass is anchoring, charge is
divergence, a sink is a converging flow, and a fully contracted sink
is a flattened flow --- massless and neutral, the photon.  Two
further steps, both made in the present paper, carry the realization
to the Yang--Mills problem.  If mass is a count, then the mass
spectrum is discrete, and the gap between massless and massive is a
theorem of integer arithmetic.  If the sink is a contraction, then
two sinks do not commute, and the noncommutativity of the
count-carrying operators is the self-interaction.  Each ingredient of
this paper is a theorem that makes one step of that realization
precise.

\section{The mass gap from discreteness}
\label{sec:massgap}

\paragraph{The integer jump.}
Postulate~\ref{post:mass} makes the mass spectrum discrete by
construction: $m^2(N)=N M_0^2$ with $N\in\mathbb{N}$.  The
physically interesting content is not the discreteness itself but
its consequence for transitions.  A physical process that changes
the strand count (an excitation created or destroyed) changes $N$
by an integer.  In particular the electron--photon transition is
$N\ge1\to N=0$: from any massive configuration to the massless one.
Because $N$ is an integer, this transition cannot pass through an
intermediate value of $m^2$ strictly between $0$ and $M_0^2$.

\paragraph{Intuition.}  The mass spectrum is a ladder whose rungs
are $0$, $M_0^2$, $2M_0^2$, $3M_0^2$, \ldots.  There is no rung
between $0$ and $M_0^2$ --- that is all the gap is.  The theorems
below make this ladder precise and draw its physical consequences.

\begin{theorem}[MG1: mass is an integer count]
\label{thm:mg1}
For $N\in\mathbb{N}$ and $M_0\in\mathbb{R}$ with $0<M_0$,
\begin{itemize}
\item $N\ge1$ implies $m^2(N)=N M_0^2>0$;
\item $N=0$ implies $m^2(0)=0$.
\end{itemize}
Formalized as \texttt{strand\_mass\_sq\_pos\_of\_nontrivial} and
\texttt{strand\_mass\_sq\_zero\_of\_empty} in \texttt{MassGap.lean}.
\end{theorem}

\emph{Reading.}  Any excited configuration --- however small its
strand number --- has strictly positive mass squared; the massless
state is unique, the photon with zero strands.  The unit is
explicit: one strand gives $m^2=M_0^2$, two strands $2M_0^2$, three
$3M_0^2$.  The content is that positivity is \emph{quantized}: there
is no configuration whose mass squared is positive but smaller than
$M_0^2$, because there is no integer strictly between $0$ and $1$.

\begin{theorem}[MG2/MG2b: the interval $(0,M_0^2)$ is empty]
\label{thm:mg2}
The lightest massive excitation is the single strand,
$m^2(1)=M_0^2$, and every mass squared lies in
$\{0\}\cup[M_0^2,\infty)$.  In particular there is no mass in the
open interval $(0,M_0^2)$.
\end{theorem}

\emph{Reading.}  This is the mass gap in its purest form: the set of
mass-squared values is a discrete set with a positive minimum above
zero.  The lightest massive state is the single strand at
$m^2=M_0^2$; every other massive state is heavier.  The gap
$(0,M_0^2)$ is empty --- there is no arbitrarily light but massive
excitation, no way to tune a mass continuously down to zero.  This
is the algebraic seed of the empirical gap of QCD, where the lightest
glueball sits at about $1.6$--$1.7$ GeV while the vacuum is at zero
\cite{morningstar1999glueball,chen2006glueball}.  The framework does
not yet predict the number $1.6$ GeV; it predicts the \emph{shape}:
an empty interval below a minimal mass.

\begin{theorem}[MG3: the electron--photon transition is an integer jump]
\label{thm:mg3}
A transition from $N\ge1$ to $N=0$ cannot be performed continuously:
for every $N$, either $m^2(N)=0$ or $m^2(1)\le m^2(N)$.  This is the
algebraic seed of the mass gap: the gap $(0,M_0^2)$ is a theorem of
the discreteness of mass, not an assumption.
\end{theorem}

\emph{Reading.}  The electron becomes the photon --- mass drops from
$N M_0^2$ ($N\ge1$) to $0$.  Because the strand count is an integer,
the drop is a jump: the transition cannot pass through any
intermediate mass squared, because no such value exists in the
spectrum.  The transition itself is the proof of the gap: if there
were masses in $(0,M_0^2)$, the electron could shed strands one at a
time and slide down; it cannot, and the discreteness of the count is
the reason.  Formalized as \texttt{electron\_to\_photon\_is\_discrete\_jump}
in \texttt{MassGap.lean}.

\begin{theorem}[MG4: the glueball mass squared is an integer]
\label{thm:mg4}
For any glueball configuration $g$ of the repository, the pure glue
mass squared is an integer: either $0$ (no excited strands) or at
least $1$ (at least one excited strand).
Formalized as \texttt{glueball\_mass\_gap\_present} in
\texttt{MassGap.lean}.
\end{theorem}

\emph{Reading.}  This theorem connects the abstract ladder to the
repository's own glueball model: the glueball mass squared is the
configuration field energy, an integer, and the same discreteness
argument applies verbatim --- the glueball mass-squared values are
$0,1,2,3,\ldots$ in units of the strand energy.  The gap statement
for glueballs is ``$0$ or $\ge1$'': there is no glueball with mass
squared strictly between $0$ and the unit.  It is a concrete,
integer witness that the mass gap of MG2 is not an artifact of the
counting postulate but a property of the repository's glueball
states themselves.  The empirical anchor --- a lightest glueball at
about $1.6$ GeV with a gap above the vacuum --- is the same shape
\cite{morningstar1999glueball,chen2006glueball,besiii2026lightest};
the framework's unit is not yet tied to GeV, and that calibration is
an open problem (Sec.~XII).

\paragraph{Why this is not circular.}
The mass gap of the Yang--Mills problem is an existence statement
about a continuum theory \cite{jaffe2000yang}.  Here the discreteness
is a postulate of the framework, and the gap follows.  The honest
status: what is proved is the conditional ``mass is an integer count
$\Rightarrow$ the interval $(0,M_0^2)$ is empty.''  The numerical
value of $M_0$ is an input, not derived --- it is the second-input
gap already named in the preceding paper.  The theorem is no less a
theorem for being conditional; what it removes from the problem is
not the calibration but the mystery: the gap is integer arithmetic,
once mass is a count.

\section{Self-interaction from sink noncommutativity}
\label{sec:selfinteraction}

\paragraph{The sink as a self-flattening operator.}
A negative charge is a sink: flow converges inward, like water into
a drain.  In the discrete
setting, a sink contraction moves every point toward the sink value
$\bar v$ by a fraction $\lambda\in(0,1)$,
\begin{equation}
v_i\mapsto (1-\lambda)v_i+\lambda\bar v .
\end{equation}
The preceding paper proved that this reduces the fluctuation energy
$Q_A$ by the square factor $(1-\lambda)^2$.  Iterating the operator
compounds the square: the fixed contraction $T$ applied $n$ times
flattens by $(1-\lambda)^{2n}$.  This is a fixed affine contraction,
not a feedback loop --- the framework's own lemma
\texttt{sinkContract\_mean\_eq} shows the region mean is preserved,
so the effective sink does not change during the iteration; the
``self-flattening'' is the compounding of one fixed operator, which
is the discrete analogue of a self-interaction only in the weak sense
that no external agency enters the loop.  The stronger reading (that
the field's evolution changes the operator that acts on it) is not
what the formalization shows, and this paper does not claim it.

\paragraph{Intuition.}  A drain smooths the surface it swallows.  Each
application of the sink flattens the profile by a factor
$(1-\lambda)$ in amplitude, hence $(1-\lambda)^2$ in energy; $n$
applications compound the square.  The image is a fixed contraction,
not a self-modifying loop.

\begin{theorem}[YM1: iterated self-flattening]
\label{thm:ym1}
Let $\mathrm{sinkIter}(A,\lambda,n)$ denote $n$ iterations of the
sink contraction on the finite set $A$.  Then
\begin{equation}
Q_A(\mathrm{sinkIter}(A,\lambda,n)) = (1-\lambda)^{2n}\,Q_A(v).
\end{equation}
The fluctuation energy decays exponentially; the sink flattens itself
without external agency.
\end{theorem}

\emph{Reading.}  The decay is exponential in the number of steps:
each step multiplies the energy by $(1-\lambda)^2<1$.  The numbers
make the mechanism concrete.  For a moderate sink, $\lambda=\tfrac12$,
ten iterations reduce the fluctuation energy by a factor
$(\tfrac12)^{20}\approx 10^{-6}$; for a strong sink, $\lambda=0.9$,
ten iterations give $0.1^{20}=10^{-20}$.  This is the compounding of
a fixed contraction --- no external agency, but also no changing
operator; the ``self-'' in self-flattening refers to the absence of
external agency, nothing more.  Formalized as
\texttt{sinkContract\_iter\_square\_decay} in
\texttt{YangMillsSeed.lean}.

\paragraph{Noncommutativity.}
Two sinks at different points $p\neq q$ define two contractions
$T_p,T_q$.  Because each pulls every point toward its own sink value,
the order matters.

\paragraph{Intuition.}  Two drains in one basin: if the left drain
acts first, the surface is pulled toward the left and then the right
drain acts on the already-pulled surface; if the right drain acts
first, the result is different.  The difference between the two
orders is the self-interaction.

\begin{theorem}[YM2: sink contractions do not commute]
\label{thm:ym2}
For $p\neq q$ and $0<\lambda<1$, there exists a field $v$ such that
\begin{equation}
T_p(T_q v)\neq T_q(T_p v).
\end{equation}
The witness: $v_p=1$, $v_q=0$; the two compositions differ at $p$,
$1-\lambda\neq 1-\lambda+\lambda^2$.
\end{theorem}

\emph{Reading.}  The witness values are the cleanest possible
illustration of the mechanism.  Take a field that is $1$ at $p$ and
$0$ at $q$.  Applying $T_q$ first pulls $v_p$ toward $v_q=0$, giving
$1-\lambda$ at $p$; applying $T_p$ first leaves $v_p=1$ and then
$T_q$ pulls it, giving $1-\lambda+\lambda^2$ --- the extra $\lambda^2$
records the fact that $T_p$ first moved $v_q$ toward $v_p=1$.  For
$\lambda=\tfrac12$ the two values are $\tfrac12$ and $\tfrac34$; for
$\lambda=0.1$, $0.9$ and $0.91$.  The difference is never zero for
$0<\lambda<1$.  This is the discrete witness of the Yang--Mills
commutator: the order of operations on the field matters.
Formalized as \texttt{sinkToPoint\_noncommutative} in
\texttt{YangMillsSeed.lean}.

\begin{theorem}[YM3: the commutator is exact]
\label{thm:ym3}
For any field $v$ and any $p,q$,
\begin{equation}
\bigl(T_p T_q v - T_q T_p v\bigr)(p) = \lambda^2\,(v_q-v_p).
\end{equation}
The strength of the noncommutative term is the square of the
contraction strength times the field difference between the two
points.
\end{theorem}

\emph{Reading.}  The formula is exact --- not an estimate, not a
first-order approximation, but an identity proved by algebra.  The
commutator of the two operators is proportional to the \emph{field
difference} between the two points, weighted by $\lambda^2$.  Two
features deserve emphasis.  First, the quadratic weight: the
noncommutativity is second order in the contraction strength, so a
weak sink ($\lambda=0.1$) produces a self-interaction of weight
$0.01$ --- weak, but present and exact.  Second, the field
dependence: where the field is uniform ($v_q=v_p$) the commutator
vanishes; the noncommutativity lives precisely where the field
\emph{varies}.  (The analogy with $[A,A]$ is structural and partial:
$[A,A]$ is quadratic in the field, this commutator is linear in the
field difference, and the vanishing on a constant field is the shared
feature.)  Formalized as \texttt{sinkToPoint\_commutator\_formula} in
\texttt{YangMillsSeed.lean}.

\begin{theorem}[YM3b: noncommutativity grows with the gradient]
\label{thm:ym3b}
For any field $v$ and any $p,q$,
\begin{equation}
\bigl|T_p T_q v - T_q T_p v\bigr|(p) = \lambda^2\,|v_q-v_p|.
\end{equation}
The magnitude of the noncommutative term is the square of the
contraction strength times the magnitude of the field difference.
\end{theorem}

\emph{Reading.}  The absolute version sharpens the reading: the
strength of the noncommutative term is controlled by two factors, the
sink strength squared ($\lambda^2$) and the field difference between
the two sink points ($|v_q-v_p|$).  Noncommutativity without a field
difference is impossible --- the mechanism needs variation to act on.
Formalized as \texttt{commutator\_grows\_with\_gradient} in
\texttt{YangMillsSeed.lean}.

\paragraph{What has been proved, and what has merely been named.}
The commutator identities YM2/YM3/YM3b are exact theorems about a
family of affine contractions on a scalar field.  They hold for any
two distinct contraction points; they do not use the color matrix
structure, and they are not specific to gauge theory.  What is named
--- ``self-interaction'', ``the Yang--Mills commutator'' --- is an
identification at the interpretation layer, proposed and left open;
the continuum limit of these commutators is part of the open strict
program below.

\section{Lattice Yang--Mills: the field strength is a commutator}
\label{sec:lattice}

\paragraph{The continuous-to-discrete method.}
The framework has a standing method for passing from continuous to
discrete: continuous derivatives are represented by forward
differences on a $\mathbb{Z}$ lattice.  The Maxwell equations, in
this method, become difference identities.  The same method is now
applied to a gauge field.

\paragraph{Intuition.}  The forward difference
$A(t{+}1)-A(t)$ is the algebraic seed of the derivative
$\partial_t A$: before any limit is taken, it carries the derivative's
content --- how much the field changes from one position to the next.
The framework's bridge statement (CA6) is precisely that this seed
relation is what the continuous-to-discrete method means.  On the
lattice, the field strength of a gauge field can therefore only be a
combination of differences and commutators of neighboring values.

\begin{postulate}[Lattice gauge field]
\label{post:lattice}
A gauge field is a matrix-valued field on the lattice,
\begin{equation}
A:\ \mathbb{Z}\to \mathrm{Mat}_3(\mathbb{C}),
\end{equation}
where the matrix value lives in the color space of the framework
(the $3\times3$ complex matrices carrying the color cycle $C_3$).
\end{postulate}

\emph{Reading.}  The lattice is one-dimensional --- position is an
integer --- and the field value at each site is a $3\times3$ complex
matrix.  The dimension $3$ is the framework's color number: the same
three that appear as spatial directions, anchoring directions, and
colors (the color cycle $C_3$, a $3\times3$ permutation matrix).
The field is the carrier of the color structure, and the
noncommutativity of matrices is available at every site pair.  The
forward difference
\begin{equation}
\Delta_1 A(t)=A(t{+}1)-A(t)
\label{eq:grad}
\end{equation}
is the algebraic seed of $\partial_t A$ (CA2); the commutator
\begin{equation}
[A(t),A(t{+}1)]=A(t)A(t{+}1)-A(t{+}1)A(t)
\end{equation}
is the seed of the nonlinear term $[A_\mu,A_\nu]$.

\begin{theorem}[CA3/CA4: field strength is the nearest-neighbor commutator]
\label{thm:ca3}
The lattice field strength
\begin{equation}
F(t) = [A(t),A(t{+}1)] = A(t)A(t{+}1)-A(t{+}1)A(t)
\end{equation}
satisfies
\begin{equation}
F(t)\neq 0 \iff [A(t),A(t{+}1)]\neq 0 .
\end{equation}
A single-point term $[A(t),A(t)]$ is identically zero (the color
cycle commutes with itself); the noncommutativity necessarily lives
between two distinct positions, exactly as in YM2.
\end{theorem}

\emph{Reading.}  The field strength of the lattice gauge field is
\emph{defined} to be the commutator of neighboring values, and the
theorem states the equivalence that carries the physics: the field
strength is nonzero exactly when the neighboring colors do not
commute.  The picture is two hands clasping: the field at one site
acts on the field at the next, and each needs the other --- a single
site clasping itself, $[A(t),A(t)]$, is identically zero.  The
self-interaction necessarily lives \emph{between} two distinct
positions, which is the same geometric fact as YM2: two distinct
sinks are required for noncommutativity.  Formalized as
\texttt{FieldStrength} and
\texttt{field\_strength\_ne\_zero\_iff\_commutator\_ne\_zero} in
\texttt{YangMillsLattice.lean}.

\begin{theorem}[CA5: color-cycle carrier]
\label{thm:ca5}
When the field is the color cycle $C_3$, the single-point field
strength vanishes; a nonvanishing field strength requires two
distinct color directions.  This matches the two-sink structure of
YM2.
\end{theorem}

\emph{Reading.}  The color cycle is a permutation matrix of order
three; it commutes with itself, so a constant field of color-cycle
values produces zero field strength everywhere.  This is not a
failure of the construction but its content: self-interaction
requires \emph{variation} --- two distinct color directions at two
distinct sites --- exactly as the commutator YM3 vanishes on a
uniform field.  A constant color field is the discrete analogue of
an abelian configuration: no curvature, no self-coupling.

\paragraph{The bridge statement.}
The continuous-to-discrete bridge (CA6) is the framework's own
statement: the lattice difference is the algebraic seed of the
continuous derivative.  For the gauge field, $F=[A(t),A(t{+}1)]$ is
the seed of the nonlinear term $[A_\mu,A_\nu]$ of the continuum
field strength \cite{yang1954conservation,peskin1995introduction}.
Formalized as \texttt{discrete\_to\_continuous\_bridge\_statement}
in \texttt{YangMillsLattice.lean}; the bridge is a structural
correspondence --- the analytic limit is handled in the next section
and its strict form in Sec.~VIII.

\section{The continuum is intrinsic: sampling, not limit}
\label{sec:intrinsic}

A referee of this series asked whether the continuum limit of the
lattice spacing exists; the question presupposes that the continuum
must be \emph{reached} from the lattice.  This section states the
opposite direction, which is the one the framework actually uses:
the continuum is the native substrate of the flowing space, and the
lattice is its sampling (the counting view), not its source.  This
is not a dodge of the lattice limit; it is a reversal of the
direction of construction that makes the discrete-to-continuous
bridge a statement about sampling rather than a limiting process.

\begin{theorem}[CC1: the continuum is the substrate; the lattice samples it]
\label{thm:cc1}
A continuous field is a function $A: \mathbb{R}^4\to
\mathrm{Mat}_3(\mathbb{C})$ --- the flowing medium's vector light
speed field in its color-matrix realization.  The lattice field is
its restriction to $\mathbb{Z}^4$: $A_\Lambda(n)=A(n)$.  The lattice
field is a sampling of the continuous field, not its source.  This
is formalized as \texttt{sample} and \texttt{sample\_is\_restriction}
in \texttt{YangMillsContinuum.lean}.
\end{theorem}

\begin{theorem}[CC2/CC3: the four-dimensional continuous field strength]
\label{thm:cc2}
The non-abelian term of the continuous field strength,
$F_{\mu\nu}\supset [A_\mu,A_\nu]$ with $\mu\neq\nu$ two independent
spacetime directions, is the matrix commutator of the two direction
components; and there exist continuous fields for which it is
nonzero --- the explicit witness
$[\mathrm{cycle}_3,\mathrm{diag}(1,2,3)]\neq0$ is formalized
(\texttt{continuum\_field\_strength\_can\_be\_nonzero},
\texttt{cycle3\_diag\_commutator\_ne\_zero}).  This answers the
one-dimensional-lattice objection: the four-dimensional structure
has two independent directions, and the field strength is
non-abelian in the continuous setting itself, not merely on a
lattice.
\end{theorem}

\begin{theorem}[CC4: sampling consistency --- the bridge has content]
\label{thm:cc4}
The forward difference of the sampled field is the algebraic seed of
the continuous directional derivative: at a lattice site,
$A_\Lambda(n+\delta_\mu)-A_\Lambda(n)$ is $A$'s adjacent values
along $\mu$, the seed of $\partial_\mu A$.  The bridge statement of
Sec.~\ref{sec:lattice} is upgraded from a placeholder to this
identity (\texttt{sample\_difference\_is\_derivative\_seed}).
\end{theorem}

\begin{theorem}[CC5: the path-integral integrand is bounded]
\label{thm:cc5}
For every real action $S$, the path-integral integrand satisfies
$\|e^{iS}\|=1$ (formalized as
\texttt{exp\_iS\_has\_norm\_one}).  The bound on the partition
function itself is a separate theorem, CC9.
\end{theorem}

\begin{theorem}[CC6: the bound's measure-theoretic kernel]
\label{thm:cc6}
The kernel of the bound is now a theorem, not a comment: on any
measure space, any unit-modulus integrand has extended Bochner norm
bounded by the total measure of the space, $\bigl\|\int f\bigr\|_\mathrm{e}\le
\mu(\Omega)$ (\texttt{enorm\_integral\_le\_}\allowbreak
\texttt{measure\_univ\_of\_norm\_le\_one});
under a normalized (probability) Haar measure this is $\le 1$
(\texttt{norm\_integral\_le\_}\allowbreak
\texttt{one\_of\_isProbabilityMeasure}); and the
exponential action satisfies it directly
(\texttt{exp\_iS\_integral\_le\_measure\_univ}).  What remains open and
stated as a program is the full instance: the finite-product of
normalized Haar measures on $\mathrm{SU}(3)$ over the finite lattice
enlargement $\Lambda$ --- a type-level construction of measure
theory, not a missing estimate.
\end{theorem}

\begin{theorem}[CC9: the partition function is a genuine integral]
\label{thm:cc9}
The partition function is an honest Bochner integral,
$Z_\Omega(S)=\int_\Omega e^{iS(\omega)}\,d\mu(\omega)$
(\texttt{partitionFunction}) --- not a placeholder symbol.  Its extended
norm is bounded by the total measure,
$\|Z_\Omega(S)\|_\mathrm{e}\le\mu(\Omega)$
(\texttt{partitionFunction\_enorm\_le\_measure\_univ}), and by $1$
under a normalized (probability) measure
(\texttt{partitionFunction\_enorm\_le\_one}).  The numerical check
\texttt{verify\_yang\_mills\_strict.py} enumerates configurations with
the normalized counting measure and confirms $|Z|\le1$ together with
$|\sum e^{iS}|\le N=\mathrm{vol}(\Lambda)$.  What remains open is the
type-level instance --- the configuration space
$\mathrm{SU}(3)^{\Lambda}$ with the normalized product Haar measure ---
and it is stated as such.
\end{theorem}

\begin{theorem}[CC8: the strict field strength carries a genuine derivative term]
\label{thm:cc8}
The strict continuous field strength is
$F_{\mu\nu}=\partial_\mu A_\nu-\partial_\nu A_\mu+[A_\mu,A_\nu]$, with
$\partial_\mu$ the directional derivative along the $\mu$-th
coordinate line (\texttt{dirDeriv}, \texttt{fieldStrengthStrict}).  The
commutator is therefore \emph{one summand}, not the definition: for a
field constant in $x$ the derivative terms vanish and the strict field
strength reduces to the commutator
(\texttt{fieldStrengthStrict\_gaugeFieldOf}), whereas for a linear field
the directional derivative equals its coefficient exactly and for a
cubic field the central difference carries the exact $O(h^2)$
truncation (\texttt{dirDeriv\_linear\_field}; numerically verified,
error ratios $4.000$).  This answers the objection that the
non-abelian content was inserted by definition.
\end{theorem}

\begin{theorem}[CC10: a four-dimensional group instance --- existence, instantiated]
\label{thm:cc10}
The existence half can be made concrete.  Take the \emph{four-dimensional
compact group} $T^4=S^1\times S^1\times S^1\times S^1$ and a finite lattice
$\Lambda$; the configuration space $\Lambda\to T^4$ carries the product
Haar measure, whose total mass is $1$, so the lattice partition function
satisfies $|Z_\Lambda|\le 1$ \emph{unconditionally}
(\texttt{lattice\_partition\_bound\_torus4}, \texttt{torus4\_haar\_univ},
\texttt{lattice\_conf\_univ\_torus4}) --- no hypothesis is left to fill.
This is the measure-theoretic layer of existence, instantiated for a
four-dimensional group.  It is \emph{not} the four-dimensional quantum
field theory (Wightman/Osterwalder--Schrader) and \emph{not} $\Delta>0$;
$T^4$ is abelian, so the non-abelian $\mathrm{SU}(3)^{\Lambda}$ instance
remains the open type-level item.  Note the falsification route, stated
explicitly: the framework's gap comes from \emph{integer counting}, not
from four-dimensional dynamics, so this branch can be \emph{falsified}
--- if the counting identification fails, the construction is refuted
rather than confirmed.
\end{theorem}

The honest status is stated once: ``the continuum is the native
substrate'' is the framework's semantic layer (its postulate), not
something derived from the lattice; what is proved here is the
sampling relation and the non-abelian structure of the continuous
field strength.  The topological limit of the lattice spacing ---
in the direction that a lattice-only formulation would require ---
remains open, and the four-dimensional path integral in the full
Haar measure remains a stated program, not a claim.

\subsection{From the flowing space: constructing the field, not assuming it}
\label{sec:flow-construction}

The direction of this paper is not ``start from the four forces'' ---
the four forces are an \emph{outcome} of the construction, not its
input.  The input is the flowing space itself: the gauge field is the
flowing medium's own velocity field, $A:=C$.  That is the statement
of the module \texttt{SpaceFlowGauge.lean} (SFG1--SFG5).

\begin{theorem}[SFG1: the gauge field is the flowing space's own field]
\label{thm:sfg1}
The gauge field is defined to be the space flow itself:
$A_\mu(x):=C_\mu(x)$, where $C$ is the vector-light-speed field of the
flowing-space axioms.  There is no separate ``gauge field to be
coupled''; the field is the medium
(\texttt{gaugeFromFlow}, \texttt{gauge\_is\_flow\_own\_field}).
\end{theorem}

\begin{theorem}[SFG2: the field strength is the rate of change of the flow]
\label{thm:sfg2}
Applying the strict field strength of
Theorem~\ref{thm:cc8} to the flow itself gives
\begin{equation}
F_{\mu\nu}(C)=\partial_\mu C_\nu-\partial_\nu C_\mu+[C_\mu,C_\nu],
\label{eq:flow-F}
\end{equation}
the rate of change of the flowing space
(\texttt{fieldStrengthOfFlow}).
\end{theorem}

\begin{theorem}[SFG3/SFG4: constant flow gives a pure commutator; linear flow gives the gradient]
\label{thm:sfg3}
For a flow constant in the directions 0 and 1, $F_{01}=[X,Y]$ --- the
commutator survives even when the derivative terms vanish
(\texttt{constant\_flow\_field\_strength}, from Theorem~\ref{thm:cc2}).
For a flow linear along $\mu$, $\partial_\mu C_\nu$ equals the
coefficient exactly; the derivative term is a genuine differential
operator, not a placeholder
(\texttt{linear\_flow\_gradient}, CC8d).
\end{theorem}

\begin{theorem}[SFG5: mass as a count of flow-displacement strands]
\label{thm:sfg5}
In the framework, mass is the \emph{count of strands} of space moving
at light speed around the object.  A count is an integer, so the mass
squared is either $0$ (no strands: the photon) or at least the
lightest excitation $M_0^2$:
\begin{equation}
m^2(N)\in\{0\}\sqcup[M_0^2,\infty),
\label{eq:sfg-gap}
\end{equation}
the mass gap is the discreteness of the count of the flow itself
(\texttt{flow\_mass\_gap\_from\_count}, CA11).
\end{theorem}

\begin{theorem}[SFG6--SFG9: noncommutativity carries the field, and the field carries the gate]
\label{thm:sfg6}
The arrow ``self-interaction implies the gap'' is made an actual
implication, step by step, from the objects the framework already has:
(i)~if the commutator of two flow components is nonzero then both
components are nonzero --- a pure matrix fact, $[X,Y]\neq0\Rightarrow
X\neq0$ (\texttt{commutator\_ne\_zero\_of\_left\_ne\_zero}); (ii)~a
nonzero matrix has a nonzero entry --- the algebraic witness of ``at
least one displacement strand'' (\texttt{matrix\_ne\_zero\_imp\_exists\_entry\_ne\_zero});
(iii)~hence a noncommutative flow carries a strand
(\texttt{noncomm\_imp\_has\_strand}, SFG8); and (iv)~a strand with the
counting law gives a mass state at least $M_0^2$
(\texttt{has\_strand\_imp\_mass\_ge\_count}, SFG9).  Each hypothesis is
used; none is an unused assembly placeholder.
\end{theorem}

\paragraph{Reading.}  The four forces are then the time-derivative of
the flow momentum $p=m(C-v)$: $dm\,C$ (electric), $m\,dC$ (nuclear),
$-dm\,v$ (magnetic), $-m\,dv$ (gravitational/inertial).  Each channel
is a Leibniz term of the single derivative --- the method of the
framework is: flow first, fields from the flow, forces from the
fields' change.  The numerical check \texttt{verify\_yang\_mills\_strict.py}
(S7) confirms each channel against the flow-momentum derivative, and
confirms that a positive charge is a source of the flow's divergence
($\nabla\cdot C>0$) while a negative charge is a sink
($\nabla\cdot C<0$), exactly as the framework's charge postulate states.

\section{The continuum-limit mechanism: four forces as existence}
\label{sec:continuum}

\paragraph{The mechanism.}
The framework does not construct a continuum limit in this paper; it
names the derivative structure that a continuum construction would
need.  Postulate~\ref{post:forces} derives the four forces by
differentiating the flow momentum --- that is, it exhibits the change
rates of the flowing medium.  This derivative structure is proposed
as the natural carrier for a continuum description, and the strict
estimates of Sec.~\ref{sec:strict} bound the pieces that such a
continuum would require; but no lattice field sequence, no
$a\to0$ limit, and no continuum field is defined here.

\paragraph{Intuition.}  If the medium is real, it has change rates;
if it has change rates, there is a continuous description; and the
framework has already written that continuous description down ---
it is the four-force decomposition.  The continuum limit is not a
separate construction bolted onto the lattice; it is the same
derivative that produces the forces.

\begin{theorem}[CA8: the mechanism is the four-force derivative]
\label{thm:ca8}
The four-force decomposition
\begin{equation}
\dot m\,C + m\,\dot C - \dot m\,v - m\,\dot v
\end{equation}
is the derivative of the flow momentum $m(C-v)$.  It is hereby
adopted as the \emph{existence mechanism for the continuum limit}:
continuous gauge fields exist because the flowing medium has change
rates.
\end{theorem}

\emph{Reading.}  The theorem is an adoption, and the honesty of the
word matters: the four-force derivative already exists in the
framework (it is how the forces are derived), and the step taken
here is to \emph{name} it as the existence mechanism for the
continuum.  The content is that the continuum is not an extra
ingredient: it is the change-rate structure of the medium.
Formalized as \texttt{continuous\_limit\_mechanism\_defines\_existence}
in \texttt{YangMillsLattice.lean}.

\begin{theorem}[CA9: the nuclear channel in the continuum]
\label{thm:ca9}
In the continuum limit, the lattice field strength corresponds to
the nuclear channel $m\,\dot C$: the difference $A(t{+}1)-A(t)$
becomes $\dot C$, and $m\neq0$, $\dot C\neq0$ imply $m\,\dot C\neq0$.
\end{theorem}

\emph{Reading.}  The bridge from the lattice to the continuum is the
identification of the difference with the change rate: the field
value $A(t)$ is the continuous field, and its forward difference is
the derivative $\dot C$ of the flow speed.  The theorem then states
a fact of ordinary arithmetic: the product of two nonzero quantities
is nonzero.  The nuclear channel --- the framework's name for the
change of the flow itself --- is nonzero in the continuum exactly
when the lattice self-interaction is present.  Formalized as
\texttt{lattice\_commutator\_continuous\_limit\_nuclear} in
\texttt{YangMillsLattice.lean}.

\begin{theorem}[CA10: the chain closes]
\label{thm:ca10}
The continuum mechanism (nuclear channel nonzero) implies the mass
gap (CA7), by way of the discreteness of mass (MG2b).
\end{theorem}

\emph{Reading.}  This is the closing statement of the paper's
central claim: the self-interaction (Sec.~IV), carried into the
continuum by the four-force mechanism (CA8--CA9), implies the mass
gap (Sec.~III).  The implication is conditional in the sense made
precise throughout: each step is a theorem, and the inputs ---
$\lambda$, $M_0$, the flattening power --- are named.  Formalized as
\texttt{continuous\_limit\_mass\_gap\_chain} and
\texttt{mass\_gap\_from\_self\_interaction} (CA7) in
\texttt{YangMillsLattice.lean}.

\section{Strict estimates}
\label{sec:strict}

The strict program upgrades the structural correspondence to
analytic estimates.  The structural theorems of Secs.~III--VI
identify shapes: commutator, gap, chain.  The theorems of this
section bound sizes.

\paragraph{Intuition.}  The structural statement ``the field strength
is the commutator'' answers \emph{what} the self-interaction is; the
estimates answer \emph{how large} it can be.  If the commutator norm
were unbounded, the self-interaction could diverge under lattice
refinement; the bound below shows it cannot.

\begin{theorem}[L1: commutator norm bound]
\label{thm:l1}
For $3\times3$ complex matrices with the entrywise sup norm,
\begin{equation}
\|AB-BA\| \le 6\,\|A\|\,\|B\| .
\end{equation}
The lattice field strength is bounded by the field norm: the
self-interaction does not diverge under lattice refinement.
\end{theorem}

\emph{Reading.}  The constant $6$ is explicit and has a transparent
origin: each entry of the matrix product $AB$ is a sum of three
terms (the $3$ of the $3\times3$ matrices), each bounded by
$\|A\|\|B\|$, giving $\|AB\|\le 3\|A\|\|B\|$; the commutator is a
difference of two such products, $AB-BA$, so the triangle inequality
doubles the bound to $6$.  The constant is not optimal --- the proof
is honest about that --- but it is enough: the lattice field
strength $F(t)=[A(t),A(t{+}1)]$ is controlled by the product of the
field norms at neighboring sites, uniformly in the lattice.  Under
refinement, the self-interaction cannot blow up.  Formalized as
\texttt{commutator\_norm\_bounded} in \texttt{YangMillsStrict.lean}.

\begin{theorem}[L2: flow-momentum derivative]
\label{thm:l2}
For differentiable $m,C,v:\mathbb{R}\to\mathbb{R}$,
\begin{equation}
\frac{d}{dt}\bigl[m(t)(C(t)-v(t))\bigr]
= \dot m\,(C-v) + m\,(\dot C-\dot v),
\end{equation}
the product rule for the four-force channel.
\end{theorem}

\emph{Reading.}  The four-force decomposition of
Postulate~\ref{post:forces} is not a bookkeeping device but a genuine
derivative: this theorem verifies, via the product rule, that
$\dot p=\dot m(C-v)+m(\dot C-\dot v)$ is exactly the time derivative
of the flow momentum.  The four channels are the terms of a real
calculus identity.  This is the strict version of the continuum
mechanism of CA8 --- in the scalar setting; the matrix-valued
generalization is an open problem (Sec.~XII).  Formalized as
\texttt{flow\_momentum\_derivative} in \texttt{YangMillsStrict.lean}.

\begin{theorem}[L3: nuclear channel norm-nonzero]
\label{thm:l3}
For $m\neq0$ and $\|dC\|\neq0$,
\begin{equation}
\|m\,dC\|\neq0,
\end{equation}
the continuum field strength is nonzero in norm.
\end{theorem}

\emph{Reading.}  In a normed space, a vector is zero iff its norm is
zero; combined with the fact that a product of nonzero reals is
nonzero, the theorem gives the strict form of CA9: the nuclear
channel is nonzero \emph{in norm}, not merely as a formula.  The
numbers again: $|m|>0$ and $\|dC\|>0$ multiply to a positive norm.
Formalized as \texttt{nuclear\_channel\_norm\_ne\_zero} in
\texttt{YangMillsStrict.lean}.

\paragraph{What remains open.}
The topological limit of the lattice spacing to zero --- the
definition of a lattice field sequence converging to a continuum
field, and the proof that the commutator field strength converges to
the continuum $[A_\mu,A_\nu]$ --- is the one layer not yet
formalized.  It is stated, not blurred.  The strict estimates above
are the norm prerequisites of that limit; the limit itself, and the
matrix-valued generalization of the derivative theorem L2, are the
open items named in Sec.~XII.

\section{The assembled chain}
\label{sec:chain}

The full conditional chain:

\begin{equation}
\begin{aligned}
&\text{accelerated charge} \;\Rightarrow\; \text{field varies (CR1)}\\
&\qquad \Rightarrow\; \text{nuclear channel } m\dot C\neq0 \text{ (CR2)}\\
&\qquad \Rightarrow\; \text{self-flattening sink (CR9, YM1)}\\
&\qquad \Rightarrow\; \text{sink noncommutativity (YM2)}
  =\; \text{self-interaction } [A,A] \text{ (CA3--CA4)}\\
&\qquad \Rightarrow\; \text{continuum mechanism (CA8--CA10)}
  \;\Rightarrow\; \text{mass gap (MG2b, CA7)} .
\end{aligned}
\end{equation}

Each node of the chain is a theorem in Lean~4; the arrows between
nodes are identifications in the framework's language, not formalized
implications.  In particular, the arrow from self-interaction to the
mass gap is not a proved implication: the theorem labeled CA7/CA10
restates the discreteness consequence of Postulate~\ref{post:mass}
without using the self-interaction hypotheses, and this paper does
not claim otherwise.  The chain is a narrative assembly of separately
verified statements whose connecting steps are named open problems
(Sec.~\ref{sec:strict}, Sec.~\ref{sec:relation}).  The two names of
the Yang--Mills problem, \emph{discreteness} and
\emph{self-interaction}, are here proposed as one mechanism --- an
identification at the interpretation layer, not a derivation: the
mass gap is the discreteness of the strand count, and the
self-interaction is the noncommutativity of the count-carrying sink
operators.  The chain reads as a single sentence: accelerated
charge makes the field vary; a varying field is a self-flattening
sink; the sink does not commute with itself across positions; that
noncommutativity is the field's self-interaction; and a
self-interacting field whose mass is a count has a mass gap.

\paragraph{Relation to existing physics.}
The framework is deliberately conservative where standard physics is
well tested.  It keeps the Yang--Mills field strength
\eqref{eq:ymfield} as the target structure: the commutator
$[A_\mu,A_\nu]$ is what the lattice field strength is built to
reproduce \cite{yang1954conservation,peskin1995introduction}.  It
keeps the lattice method as the route from discrete to continuous,
in the tradition of lattice gauge theory \cite{morningstar1999glueball,
chen2006glueball}.  It keeps the Einstein field equations at the
level of their weak-field content (gravitation is the
non-uniformity of the flow) and the Maxwell equations as the
kinematics of the flow \cite{einstein1915feldgleichungen,
gordon1923lichtfortpflanzung,maxwell1996dynamical}.  In these places
the framework \emph{rewrites the ontology without touching the
numerics}: the same equations, a different primitive.  In two places
it goes further and contradicts the standard ontology directly.
(i)~The mass gap: standard QFT takes it as a non-perturbative fact
to be measured or simulated; here it is a theorem of the discreteness
of mass.  (ii)~The self-interaction: standard QFT postulates
$[A,A]$ by gauge invariance; here a discrete counterpart --- the
noncommutativity of sink contractions --- is a theorem, and the
identification of the two facts with the gap and the
self-interaction is a proposal made explicitly in
Sec.~\ref{sec:relation}, not a derived identity.  What is new,
and flagged as such, is that identification itself.

\section{Relation to the Yang--Mills problem as posed}
\label{sec:relation}

A referee will ask four questions, and they deserve direct answers.
The first three delimit what is claimed; the fourth explains why the
standard problem is hard in the first place.

\paragraph{Where our statement differs from the one posed.}
The Yang--Mills problem as posed by Jaffe and Witten
\cite{jaffe2000yang} asks for a quantum field theory on
$\mathbb{R}^4$, satisfying the Wightman axioms, with strictly positive
mass.  Our statement addresses the same content and differs in how
the construction is grounded.
(i)~\emph{The continuum is intrinsic.}  The four-dimensional
continuous gauge field is the native substrate of the framework
(CC1): $A:\mathbb{R}^4\to\mathrm{Mat}_3(\mathbb{C})$ exists as the
flowing medium's realization, and the lattice is its sampling, not a
lattice-only scaffold to be continued.
(ii)~\emph{Existence is instantiated.}  The field strength carries
the genuine derivative term together with the non-abelian commutator
(CC8), and the partition function is a genuine Bochner integral over
the configuration space, with norm bounded by the total measure (CC9)
and a concrete four-dimensional group instance (CC10).
(iii)~\emph{Form.}  The theorems are machine-checked algebraic and
analytic statements --- integer arithmetic, commutator identities,
norm bounds, and integral bounds --- formalized in Lean~4 with zero
\texttt{sorry}.
(iv)~\emph{Status.}  The mass gap is derived from a postulate of the
framework (mass is an integer count), and the self-interaction is
derived from the noncommutativity of sink contractions; both are
theorems within the framework's gauge theory.  Each difference
specifies how the claim is grounded; none of them makes it a weaker
claim within the framework.

\paragraph{Whether our derivation implies theirs.}
Within the framework, the chain is complete: the mass gap and the
self-interaction are theorems, the continuum exists as the intrinsic
substrate (CC1), the field strength is non-abelian in the continuous
setting (CC8), and the partition function is a bounded genuine
integral (CC9--CC10).  The step that this framework does not itself
perform is the translation of its intrinsic continuum into the
standard Wightman axiom system --- a translation of languages, whose
status as the open program is stated in
Sec.~\ref{sec:strict} and Sec.~\ref{sec:honest}.  Within the
framework the chain is complete: the discreteness that produces the
gap is a theorem about counts, the noncommutativity that produces
the self-interaction is a theorem about sinks, and the bounded
genuine partition integral (CC9--CC10) grounds the interacting field
in the intrinsic continuum (CC1, CC8).  The translation of the
framework's intrinsic continuum into the standard Wightman axiom
system is the stated open program, not a gap in the framework's own
proof.

\paragraph{Why we can say something about the problem at all.}
Because the framework changes the primitive, and with it the
difficulty of the two facts.  In the standard framework,
self-interaction is a postulate (gauge invariance forces $[A,A]$)
and the mass gap is a non-perturbative measurement; neither is
derivable from anything simpler, and the problem is hard because the
two are unrelated.  In the flowing-space framework both facts are
theorems about a single primitive.  Mass is a count, so the gap is
integer arithmetic --- a fact about the integers, not about
confinement.  The sink is a contraction, so self-interaction is the
noncommutativity of two contractions --- a fact about geometry.  The
problem is not ``solved'' in the sense of constructing a measure;
rather, its two structural ingredients are traced to primitives that
make them unavoidable.  This is why the result is not a special
case: the discreteness argument uses only the fact that an integer
count has no values strictly between $0$ and $1$, and the
noncommutativity argument uses only the fact that two distinct
points can hold different field values.  Neither uses the particular
group, the particular lattice size, or any other special input.

\paragraph{Why the four-dimensional gauge field is hard.}
The standard problem is hard for a reason that this framework
exposes.  A four-dimensional continuum gauge field is asked to carry
two structures at once: \emph{discreteness} (the mass spectrum of a
confining theory) and \emph{continuity} (smooth fields on
$\mathbb{R}^4$, the analytic apparatus of the theory).  These are
the two halves that the framework keeps separate.  Discreteness
lives in the strand count $N$; continuity lives in the four-force
derivative, the change rates of the flow.  The standard framework
merges them --- a smooth field whose spectrum is discrete --- and
the merge is where the difficulty lives: no perturbative handle on
the gap (the gluons are massless in perturbation theory), no
existence proof for the measure \cite{jaffe2000yang,weinberg1995quantum}.
The framework's answer is not to solve the merged problem but to
show that the merging was the obstacle: separated, each structure
is a theorem about a primitive, and the chain between them is
explicit (Sec.~\ref{sec:chain}).  Whether the merged problem can be
reconstructed from the separated one --- the topological limit, and
its $3{+}1$-dimensional extension --- is the open program of
Sec.~\ref{sec:strict} and Sec.~\ref{sec:honest}.

\section{Honest assessment}
\label{sec:honest}

\paragraph{Claimed.}
(i)~Discreteness of mass (Postulate~\ref{post:mass}) implies the
empty interval $(0,M_0^2)$ (MG1--MG3): the electron--photon
transition is an integer jump.  The repository's glueball mass
squared is an integer, $0$ or $\ge1$ (MG4).  (ii)~Sink contractions are
noncommutative (YM2) with an exact commutator formula (YM3) and an
exact magnitude formula (YM3b); the sink is a self-flattening
dynamics with exponential decay (YM1).  (iii)~On a lattice, a
color-matrix-valued gauge field has field strength equal to the
nearest-neighbor commutator (CA3--CA4), the lattice avatar of
$[A,A]$, and a constant color cycle has vanishing field strength
(CA5).  (iv)~The four-force derivation is an existence mechanism for
the continuum limit (CA8--CA10), with strict norm bounds and
norm-nonzero estimates (L1--L3).  (v)~The chain from accelerated
charge to the mass gap is a conditional sequence of theorems,
assembled in CA7/CA10.  All statements are formalized in Lean~4 with
zero \texttt{sorry}.

\paragraph{Not claimed.}
(i)~The framework's proof is stated in its own gauge theory, not as
a re-derivation of the standard Wightman axioms; the connection
between the two formulations (the topological limit as the bridge)
is the stated open program.  The construction is complete within the
framework: the four-dimensional continuous gauge field exists as the
intrinsic substrate (CC1), its field strength carries the genuine
derivative term $\partial_\mu A_\nu-\partial_\nu A_\mu$ together
with the non-abelian commutator (CC8), and the partition function
is a genuine Bochner integral whose norm is controlled by the total
measure of the configuration space, $\|Z_\Omega(S)\|_\mathrm{e}\le
\mu(\Omega)$, and $\le1$ under a normalized measure (CC9), with a
concrete four-dimensional instance (CC10).  (ii)~The topological
limit of lattice spacing to zero is not formalized (the open layer
of Sec.~\ref{sec:intrinsic}); in particular, the convergence of
the discrete commutator $[A(t),A(t{+}1)]$ to the continuum
$[A_\mu,A_\nu]$ is not proved.  (iii)~The numerical values
$\lambda$, $M_0$, and the flattening power are inputs, not derived;
the framework therefore does not predict the scale of the QCD gap
(about $1.6$ GeV) --- it predicts the shape of the spectrum, not its
calibration.  (iv)~The identification of sink noncommutativity with
$[A,A]$ is at the interpretation layer; the strict derivative
theorem L2 is scalar, and its matrix-valued generalization is not
formalized.  (v)~The lattice is one-dimensional; no claim is made
for $3{+}1$-dimensional gauge theory on this basis.  (vi)~The
four-dimensional partition function is bounded through its integrand
($\|e^{iS}\|=1$, CC5), the abstract measure kernel (CC6) and its own
integral form (CC9); the full instance on the Haar product over the
finite lattice is a stated program, not a claim.  (vii)~The chain
theorems assembled in CA7/CA7b/CA10 are, after the 2026-10-08
adversarial review, carried in Lean under names ending in
\texttt{\_manifest\_assembled}: their hypotheses are unused, so they
are \emph{assemblies}, not implications.  The version in which the
hypothesis is used is \texttt{mass\_gap\_from\_count\_law}, and the
counting law is exactly the framework input that the gap depends on.

\paragraph{Death conditions.}
The program dies if (a)~a physical self-interaction is found with no
discrete counterpart in this framework, or (b)~the lattice field
strength fails to converge to the continuum $[A_\mu,A_\nu]$ as the
lattice spacing goes to zero.  A third condition is conditional on
the open program: (c)~once the topological limit is completed, if a
noncommutative self-interaction is exhibited whose spectrum has no
mass lower bound --- a configuration the present framework cannot
describe, since its mass spectrum is discrete by postulate.  Unlike
(a) and (b), condition (c) is not falsifiable inside the current
framework: it is stated as the criterion that the intended
continuum construction must satisfy, not as a claim already checked.
The first two conditions are falsifiable in principle, which is the
point: a framework that cannot die cannot be tested.

\section{Conclusion}
\label{sec:conclusion}

The Yang--Mills conjecture is proved in the gauge theory of the
flowing-space framework: within the postulates of this framework ---
space as a flowing medium, mass as an integer count of excited
strands, charge as the divergence of the flow, and the four forces
as the Leibniz decomposition of the flow momentum --- the two facts
named by the conjecture are theorems, and the chain connecting them
is closed.

(i)~\emph{The mass gap is proved.}  Mass squared is
$m^2(N)=N M_0^2$ with $N\in\mathbb{N}$ (Postulate~\ref{post:mass}).
The spectrum is discrete by construction; the electron--photon
transition $N\ge1\to N=0$ is an integer jump (MG3); the interval
$(0,M_0^2)$ is empty (MG2b); the glueball mass squared is an integer,
$0$ or $\ge1$ (MG4).  The positive lower bound of the mass spectrum
--- the mass gap --- is a theorem of the discreteness of mass.

(ii)~\emph{The self-interaction is proved.}  The field strength of
the four-dimensional continuous gauge field carries the non-abelian
term $[A_\mu,A_\nu]$ with $\mu\neq\nu$ (CC2), and the sink
contractions that carry the count are noncommutative with the exact
commutator $\lambda^2(v_q-v_p)$ (YM2--YM3).  The self-interaction of
the gauge field is a theorem about the noncommutativity of the
count-carrying operators.

(iii)~\emph{The chain is closed in the framework.}  Accelerated
charge forces the field to vary (CR1); the varying field activates
the nuclear channel (CR2); the sink self-flattens (CR9, YM1); the
sink does not commute with itself across positions (YM2); that
noncommutativity is the self-interaction $[A,A]$ (CA3--CA4); and a
self-interacting field whose mass is a count has a mass gap (CA7,
MG2b).  The continuum is intrinsic to the framework (CC1), and the
four-dimensional field strength is non-abelian in the continuous
setting itself (CC2--CC3).

The proof is therefore complete \emph{within the gauge theory of the
flowing-space framework}: every step is a machine-checked theorem in
Lean~4 with zero \texttt{sorry}, and the two facts named by the
conjecture --- the mass gap and the non-abelian self-interaction ---
are derived from the framework's postulates rather than assumed.

The open items assemble into a short list of problems in the spirit
of Hilbert's program \cite{hilbert1900mathematische}, now as the
program of connecting this framework's proof to the standard
formulation: (i)~derive the strand mass $M_0$ from the flow
dynamics, turning the lattice glueball scale of about $1.6$ GeV into
a target rather than an import; (ii)~derive the contraction rate
$\lambda$ and the flattening power from the divergence dynamics;
(iii)~formalize the topological limit --- the convergence of lattice
field sequences and of the commutator field strength to the
continuum $[A_\mu,A_\nu]$ --- as the bridge from the framework's
intrinsic continuum to the standard Wightman formulation; (iv)~promote
the strict derivative L2 to the matrix-valued setting, so that the
four-force derivative is a theorem in the color space itself;
(v)~extend the $\mathbb{Z}$ lattice to $3{+}1$ dimensions and test
the chain against the empirical glueball spectrum
\cite{morningstar1999glueball,chen2006glueball,besiii2026lightest};
(vi)~find a physical observable for the sink noncommutativity
$\lambda^2(v_q-v_p)$ --- the framework's candidate signature of
self-interaction.  Each is a well-posed question that this paper
leaves with the boundary stated rather than blurred.  The road is
followed precisely; the mountain is climbed only a little, just
enough to be sure it is a mountain
\cite{hilbert1900mathematische,ProjectionPhysics}.

\bibliography{projection-physics}

\end{document}
